%% ma: L12-B08-0148
%% lop: 12
%% chuong: 2
%% bai: 8
%% dang: 12.8.50
%% loai: TN4
%% muc: VD
%% dap_an: "D"
%% nguon: "CĐ2.3.pdf giáo viên gửi 27/10/2026 (M07, Lê Minh Khôi – Nguyễn Anh Khôi, Chương II Chủ đề 3 Biểu thức tọa độ của các phép toán vectơ), Câu 64"
%% kien_thuc: "Bài 8 (điểm chia đoạn thẳng, độ dài), lớp 10 (tính chất đường phân giác)"
%% trang_thai: cho-duyet
%% ly_do: "Dạng mới 12.8.50 đề xuất, chờ giáo viên duyệt"
\begin{ex}
Trong không gian $Oxyz$, cho tam giác $ABC$ có tọa độ các đỉnh $A(1;1;1)$, $B(5;1;-2)$, $C(7;9;1)$. Tính độ dài đường phân giác trong $AD$ của góc $A$.
\choice{$AD=\dfrac{3\sqrt{74}}2$}{$AD=2\sqrt{74}$}{$AD=3\sqrt{74}$}{\True $AD=\dfrac{2\sqrt{74}}3$}
\loigiai{$AB=5$, $AC=10$ nên $\dfrac{DB}{DC}=\dfrac12\Rightarrow\overrightarrow{DC}=-2\overrightarrow{DB}\Rightarrow D=\dfrac{2B+C}3=\left(\dfrac{17}3;\dfrac{11}3;-1\right)$.

$\overrightarrow{AD}=\left(\dfrac{14}3;\dfrac83;-2\right)$, $AD=\dfrac{\sqrt{296}}3=\dfrac{2\sqrt{74}}3$.}
\end{ex}
